% Pursuing Stacks (cotes 134-2 to 134-8): archive cote 134-7, site page 90 (PDF page 91), typescript page 532.
% typescript. Draft compared: ps6.tex:4533-4590.
% AI-transcribed by Claude Opus 5.5 (claude-opus-5-5, adaptive thinking, effort high), Message Batches API, batch ps4 (msgbatch_012Aoc1NuskATPsyBeKSYK7J), 2026-11-12; not yet checked by eye.
\scanpage{134-7/90}
\mspage{532}

A admits enough projectives or enough injectives respectively,  these
categories\inserted[typed]{(as pointed out by Quillen)} are associated to closed model structures on $C^{-}(A)$ and $C^{+}(A)$
respectively,  \deleted{and} \inserted{but (it seems)} not otherwise.

It strikes me that nobody apparently so far has tried to develop
homotopy theory,  starting as basic data with a category of models \xout{a} $M$
together with a set $W$ of weak equivalences in $M$ ,  satisfying suitable
assumptions,  without giving moreover a notion of ``fibration'' or of
``cofibrations'' in M.  Various examples,   indluding the all-important
case (from my point of view) of (Cat),  suggest that the choice of the
notion of fibration or cofibration is'nt really so imperative,  that it
is to a certain extent arbitrary,  different choices (compatible with the
basic data W) leading to essentially\hand{/}``the same''\hand{/}homotopy theory.  This \xout{xxxxx}
is due to the fact that they lead to the same \unsure{notions} of ``integration'' and
``cointegration'' of homotopy types,  which \unsure{depend\xout{s}|depends} indeed only on W and
which,  in my eyes,  \deleted{\xout{is}}\inserted{are} \uline{the} two main operations of homotopy theory
(compare section 69).  They seem to me the key for defining such a thing
as a \inserted[typed]{specific} ``homotopy theory'' ,  independently of any particular choice of
a category of models ( + extra structure\inserted[typed]{on it}, and notably\hand{,}weak equivalences)
used for describing it.  The precise technical notion achieving this is
of course the notion of a ``derivator'' - and I do hope that it should'nt
be too awfully hard to construct,  for instance,  the ``derivator of
schematic homotopy types'',  and maybe even \unsure{\over{ch}{ch}}aracterize it axiomatically
(as well as the derivator of usual homotopy types),  without having to
\inserted[typed]{make any mention of models}\deleted{\xout{xxxx}} for such characterization.  One may say that,  after \xout{the} \inserted{the}\inserted[typed]{one} major
step in \inserted[typed]{the} foundations of homological algebra,\hand{|}consisting in introducing \xout{and}
the derived category of an abelian category (and systematically working
with \xout{this} derived categories for stating the main \xout{theorems} facts about
cohomo\unsure{\hand{lo}}gy of all kind of ``spaces'' namely topoi,  such as usual topologica\unsure{l}
spaces,  schemes and the like ...),  the \xout{foundational} work on foundations
more or less stopped short,  while the next step to take was to come to
a grasp of the full structure involved in derived categories,  namely the
structure of a derivator.  And it turns out that \xout{both}\inserted[typed]{as well} the step of taking
derived categories (namely localizations of suitable model categories),
as the next namely introducing the corresponding structure of a derivator\hand{,}
make good sense also in non abelian set-ups,  namely for doing \xout{homo}
``homotopy theory'',  thus viewed as a non-comm\unsure{\hand{u}}tative  version of homology
theory.

\margin{90}
